%% Interstitial: Paper 0 -> Paper 1
%% Philosophical and conceptual bridge: probability to prediction.

Chapter~\ref{chap:paper-zero} established that compatible observable
laws determine a unique $\sigma$-additive probability measure $P$ on
$(\Omega, \sigma(\mathcal{Q}))$.  The measure exists.  It is canonical.
Its marginals recover the original valuations $\{\ell_Q\}$, and two
measures agreeing on all query marginals are identical: there is no
hidden probability beyond what the observable interface shows.

\medskip

\paragraph{The residual freedom.}
What $P$ does not determine on its own is the \emph{temporal structure}
of the observer's queries.  The observer measures at a time and the
system evolves; the canonical measure $P$ is the probability over the
realization space $\Omega$, but $\Omega$ encodes outcomes, not futures.
Given the current query $Q$ and some future observable $F$, the measure
$P$ determines a conditional distribution of $F$ given $Q$.  But this
conditional distribution depends on the full measure $P$ and the entire
query system.  The predictive content of $Q$ --- what $Q$ alone can
tell the observer about $F$ --- appears to require a choice: a choice
of which aspect of $P$ to attend to, and through what minimal lens.

The observer, it seems, retains a freedom: how their present measurement
relates to future expectation.  Whether there is a minimal such structure
--- one that is not chosen but read off from $Q$ and $P$ together ---
is not settled by the existence of $P$.

\medskip

\paragraph{The elimination.}
Chapter~\ref{chap:paper-one} shows this freedom is illusory.  The
predictive content of $Q$ is entirely fixed by $(Q, P)$.  The
\emph{minimal predictive query} $Q_* = \varphi_Q \circ Q$ is the
coarsest function of $Q$ that retains all of $Q$'s predictive
information about any future observable: it is not constructed by
the observer but derived as the canonical quotient of $Q$ under
predictive equivalence.  The \emph{predictive factorization theorem}
confirms that every bounded measurable $g(F)$ satisfies
\[
  \mathbb{E}[g(F) \mid \sigma(Q)] = (\mathrm{predictiveOp}\,g) \circ Q_*,
\]
and the factorization is unique through $Q_*$.  The observer's
present-to-future orientation is not a free act; it is the form
their present measurement takes under the constraint of coherent
probability.

\medskip

\paragraph{The structure this introduces.}
This is the first result in the program that is genuinely dynamical in
character.  The passage from observable compatibility to predictive
sufficiency is the passage from static coexistence of queries to the
temporal structure of prediction.  The query $Q$ represents the observer's
present measurement; $Q_*$ is the minimal structure that makes the future
legible from the present; the predictive kernel $\Pi_Q$ encodes what is
visible through that lens.

What remains undetermined after Paper~1 is not the predictive structure
of any single query, but whether the predictive structures at different
times cohere --- whether there is a consistent temporal \emph{dynamics}
relating $Q_*$ at time $0$ to the predictive content at time $t$.
That is the question Chapter~\ref{chap:paper-two} takes up.
